\documentclass[10pt,a4paper]{amsart}
\usepackage{amsmath,amssymb,amsthm,graphicx}
% Retained minimal preamble for the Stacks Project chapter varieties.tex.
% Verbatim excerpt of the chapter preamble: Xy-pic setup (native preamble lines 14-18), the native theorem
% declarations with the mathematical macros (native preamble lines 154-217), and the project own no-op
% definitions of the reference/slogan/history display environments (native wrapper.tex lines 1-3) that the
% counted statement uses for its slogan. Copied mechanically; nothing added, removed or redefined.
% For commutative diagrams we use Xy-pic
\usepackage[all]{xy}

% We use 2cell for 2-commutative diagrams.
\xyoption{2cell}
% Theorem environments.
%
\theoremstyle{plain}
\newtheorem{theorem}[subsection]{Theorem}
\newtheorem{proposition}[subsection]{Proposition}
\newtheorem{lemma}[subsection]{Lemma}

\theoremstyle{definition}
\newtheorem{definition}[subsection]{Definition}
\newtheorem{example}[subsection]{Example}
\newtheorem{exercise}[subsection]{Exercise}
\newtheorem{situation}[subsection]{Situation}

\theoremstyle{remark}
\newtheorem{remark}[subsection]{Remark}
\newtheorem{remarks}[subsection]{Remarks}

\numberwithin{equation}{subsection}

% Macros
%
\def\lim{\mathop{\mathrm{lim}}\nolimits}
\def\colim{\mathop{\mathrm{colim}}\nolimits}
\def\Spec{\mathop{\mathrm{Spec}}}
\def\Hom{\mathop{\mathrm{Hom}}\nolimits}
\def\Ext{\mathop{\mathrm{Ext}}\nolimits}
\def\SheafHom{\mathop{\mathcal{H}\!\mathit{om}}\nolimits}
\def\SheafExt{\mathop{\mathcal{E}\!\mathit{xt}}\nolimits}
\def\Sch{\mathit{Sch}}
\def\Mor{\mathop{\mathrm{Mor}}\nolimits}
\def\Ob{\mathop{\mathrm{Ob}}\nolimits}
\def\Sh{\mathop{\mathit{Sh}}\nolimits}
\def\NL{\mathop{N\!L}\nolimits}
\def\CH{\mathop{\mathrm{CH}}\nolimits}
\def\proetale{{pro\text{-}\acute{e}tale}}
\def\etale{{\acute{e}tale}}
\def\QCoh{\mathit{QCoh}}
\def\Ker{\mathop{\mathrm{Ker}}}
\def\Im{\mathop{\mathrm{Im}}}
\def\Coker{\mathop{\mathrm{Coker}}}
\def\Coim{\mathop{\mathrm{Coim}}}

% Boxtimes
%
\DeclareMathSymbol{\boxtimes}{\mathbin}{AMSa}{"02}

%
% Macros for moduli stacks/spaces
%
\def\QCohstack{\mathcal{QC}\!\mathit{oh}}
\def\Cohstack{\mathcal{C}\!\mathit{oh}}
\def\Spacesstack{\mathcal{S}\!\mathit{paces}}
\def\Quotfunctor{\mathrm{Quot}}
\def\Hilbfunctor{\mathrm{Hilb}}
\def\Curvesstack{\mathcal{C}\!\mathit{urves}}
\def\Polarizedstack{\mathcal{P}\!\mathit{olarized}}
\def\Complexesstack{\mathcal{C}\!\mathit{omplexes}}
% \Pic is the operator that assigns to X its picard group, usage \Pic(X)
% \Picardstack_{X/B} denotes the Picard stack of X over B
% \Picardfunctor_{X/B} denotes the Picard functor of X over B
\def\Pic{\mathop{\mathrm{Pic}}\nolimits}
\def\Picardstack{\mathcal{P}\!\mathit{ic}}
\def\Picardfunctor{\mathrm{Pic}}
\def\Deformationcategory{\mathcal{D}\!\mathit{ef}}
\newenvironment{reference}{}{}
\newenvironment{slogan}{}{}
\newenvironment{history}{}{}
\usepackage{cleveref}
\crefname{originallemma}{original lemma}{original lemmas}
\Crefname{originallemma}{Original lemma}{Original lemmas}
\makeatletter
\newlabel{algebra-lemma-flat-under-regular}{{00OF}{source}{}{sourcepointer-1}{}}
\newlabel{algebra-lemma-flat-under-regular@cref}{{[originallemma][0][]00OF}{[1][0][]source}{}{}{}}
\makeatother
\crefname{originallemma}{original lemma}{original lemmas}
\Crefname{originallemma}{Original lemma}{Original lemmas}
\makeatletter
\newlabel{lemma-geometrically-regular}{{038V}{source}{}{sourcepointer-2}{}}
\newlabel{lemma-geometrically-regular@cref}{{[originallemma][0][]038V}{[1][0][]source}{}{}{}}
\makeatother
\crefname{originallemma}{original lemma}{original lemmas}
\Crefname{originallemma}{Original lemma}{Original lemmas}
\makeatletter
\newlabel{morphisms-lemma-base-change-module-flat}{{01U8}{source}{}{sourcepointer-3}{}}
\newlabel{morphisms-lemma-base-change-module-flat@cref}{{[originallemma][0][]01U8}{[1][0][]source}{}{}{}}
\makeatother
\title{Geometric regularity descends along a flat morphism of locally Noetherian schemes over a field}
\author{The Stacks Project Authors}
\date{stacks-project:varieties.tex:local-snapshot-20261011; GNU FDL 1.2 or later}
\begin{document}
\maketitle
\noindent\textit{Editorial note.} All delivered mathematical text is the author's own, copied from the pinned Stacks Project chapter file varieties.tex; no statement or proof step was rewritten, reordered or removed.\par\medskip
\noindent\textit{Editorial note.} The chapter text immediately before the lemma notes that it is a geometric variant of an algebra lemma; that remark lies outside the counted span and enters only as context.\par\medskip
\noindent\textit{Editorial note.} The counted spans carry 3 navigation stubs and no bibliography citation.\par\medskip
\noindent\textit{Editorial note.} The preamble is the chapter's retained minimal preamble (sources/stacks-varieties/retained-preamble.tex lines 1-77), a verbatim excerpt of the native preamble.\par\medskip
\noindent\textit{Editorial note.} The label lemma-flat-under-geometrically-regular is defined only in varieties.tex, so no cross-chapter restatement had to be weighed for this card.\par\medskip
\noindent\textit{Allowed standard prerequisites (editorial).}\begin{itemize}
\item Regularity descends along flat local homomorphisms of Noetherian local rings. Algebra, Lemma algebra-lemma-flat-under-regular, invoked at proof lines 2537-2540. Supplies the descent of regularity on local rings.
\item The pointwise criterion for geometric regularity. Same chapter, Lemma lemma-geometrically-regular, invoked at proof lines 2525-2527. Supplies the passage from regularity after all finite purely inseparable base changes to geometric regularity.
\item Flatness is preserved by base change. Morphisms, Lemma morphisms-lemma-base-change-module-flat, invoked at proof lines 2528-2530. Supplies that the base changed morphism is still flat at the lifted point.
\end{itemize}\medskip
\ifdefined\hypertarget\hypertarget{sourcepointer-1}{}\fi
\noindent\textit{Source reference.} Original lemma 00OF: sources/stacks-chow/tags.txt, source lines 893--893. Pointer outside the counted span, invoked once for context; the load-bearing argument of the counted statement is the author's proof reproduced in full in the proof span.\par\medskip
\ifdefined\hypertarget\hypertarget{sourcepointer-2}{}\fi
\noindent\textit{Source reference.} Original lemma 038V: sources/stacks-varieties/pinned-original-source.tex, source lines 2424--2424. Pointer outside the counted span, invoked once for context; the load-bearing argument of the counted statement is the author's proof reproduced in full in the proof span.\par\medskip
\ifdefined\hypertarget\hypertarget{sourcepointer-3}{}\fi
\noindent\textit{Source reference.} Original lemma 01U8: sources/stacks-chow/tags.txt, source lines 2401--2401. Pointer outside the counted span, invoked once for context; the load-bearing argument of the counted statement is the author's proof reproduced in full in the proof span.\par\medskip
\begin{lemma}
\label{lemma-flat-under-geometrically-regular}
Let $k$ be a field.
Let $f : X \to Y$ be a morphism of locally Noetherian schemes over $k$.
Let $x \in X$ be a point and set $y = f(x)$.
If $X$ is geometrically regular at $x$ and
$f$ is flat at $x$ then $Y$ is geometrically regular at $y$.
In particular, if $X$ is geometrically regular over $k$ and
$f$ is flat and surjective, then $Y$ is geometrically regular over $k$.
\end{lemma}

\begin{proof}
Let $k'$ be finite purely inseparable extension of $k$.
Let $f' : X_{k'} \to Y_{k'}$ be the base change of $f$.
Let $x' \in X_{k'}$ be the unique point lying over $x$.
If we show that $Y_{k'}$ is regular at $y' = f'(x')$, then
$Y$ is geometrically regular over $k$ at $y'$, see
Lemma \ref{lemma-geometrically-regular}.
By
Morphisms, Lemma \ref{morphisms-lemma-base-change-module-flat}
the morphism $X_{k'} \to Y_{k'}$ is flat at $x'$.
Hence the ring map
$$
\mathcal{O}_{Y_{k'}, y'}
\longrightarrow
\mathcal{O}_{X_{k'}, x'}
$$
is a flat local homomorphism of local Noetherian rings with
right hand side regular by assumption. Hence the left hand side
is a regular local ring by
Algebra, Lemma \ref{algebra-lemma-flat-under-regular}.
\end{proof}

\end{document}
